%%%%%%%%%%%%%%%%%%%%%%%%%%%%%%%%%%%%%%%%%%%%%%%%%
\section{Knowledge Distillation and Embedded Inference Feasibility}\label{sec:lit_kd_and_embedded}

Knowledge distillation and embedded-hardware feasibility were originally planned as two substantial, semi-independent research threads: a survey spanning feature-based detection-specific distillation, capacity-gap mitigation via teacher-assistant chains, multi-teacher and cross-domain distillation strategies, and a full quantization-aware-training pass targeting the \textit{QCS605}'s Hexagon 685 DSP. With the majority of the available time consumed by dataset construction and synthetic-data generation (\Cref{sec:lit_synthetic_data}), what remained was room for only a single, basic teacher-student distillation comparison and a single runtime feasibility check on real target-adjacent hardware -- reported together as one small, explicitly-scoped-as-basic secondary experiment in \Cref{sec:results_kd_and_embedded} rather than as two independent research contributions. This section accordingly restricts itself to the foundational material needed to justify and interpret those two basic experiments, rather than surveying the wider design space that motivated the original, more ambitious plan.

%%%%%%%%%%%%%%%%%%%%%%%%%%%%%%%%%%%%%%%%%%%%%%%%%
\subsection{Knowledge Distillation Fundamentals}\label{sec:lit_knowledge_distillation}

Knowledge distillation traces to an early model-compression technique in which a compact classifier is trained to reproduce the predictions of a large, slow ensemble rather than the original hard labels \cite{buciluaModelCompression2006}, and was reformulated into its now-standard temperature-scaled soft-target form by Hinton, Vinyals, and Dean \cite{hintonDistillingKnowledgeNeural2015}. Raising the temperature $T$ of the teacher's softmax before computing the student's training targets exposes the teacher's \enquote{dark knowledge}: the relative probabilities it assigns to \emph{incorrect} classes, which encode a similarity structure over the label space that a one-hot ground-truth label cannot (an image of a wolf may have a small probability of being mistaken for a dog, but that mistake is still far more probable than mistaking it for a car). The student is trained on a weighted combination of two cross-entropy losses -- one against the teacher's softened targets, computed at the same high temperature, and one against the ground-truth labels at a temperature of one -- with the soft-target term's gradient rescaled by $T^2$ so the two terms remain comparable as $T$ is varied.

Subsequent work organizes distillation strategies by which part of the teacher is used as supervision, a taxonomy systematized in the survey by Gou, Yu, Maybank, and Tao \cite{gouKnowledgeDistillationSurvey2021}: \emph{response-based} knowledge, which -- as in Hinton et al.'s original formulation -- matches only the teacher's final-layer outputs; \emph{feature-based} knowledge, which additionally supervises the student's intermediate representations; and \emph{relation-based} knowledge, which matches structural relations between samples or feature-map pairs rather than individual predictions. Feature-based distillation was introduced by \textit{FitNets} \cite{romeroFitNetsHintsThin2015}, in which a (possibly convolutional) regressor maps a chosen intermediate \enquote{guided} layer of the student onto a \enquote{hint} layer of the teacher; because this hint-based pretraining stage gives the student a better parameter initialization than random weights before response-based distillation is applied, it allows training students that are deeper and thinner than the teacher, not merely smaller, at moderate cost in additional design choices (which layer pair to match, and how to project between mismatched feature-map shapes).

Adapting response-based distillation to object detection is not a straightforward transplant of the classification recipe. A detector's output combines dense per-location classification and regression predictions under extreme foreground-background imbalance, and distilling teacher and student features uniformly across this imbalance is empirically counterproductive: Focal and Global Distillation (FGD) \cite{yangFocalGlobalKnowledge2022} shows that distilling foreground and background regions together, with equal weighting, performs \emph{worse} than distilling either region on its own, because the mismatch between teacher and student attention is concentrated disproportionately in the foreground. Zhendong Yang et al.'s FGD instead separates the feature map into foreground and background components using the ground-truth boxes, applies teacher-derived spatial and channel attention masks so that the student's limited capacity is concentrated on the teacher's most salient pixels and channels, and adds a global relational term (built from pixel-pair correlations) to recover the cross-region context that a hard foreground/background split otherwise discards. This line of work establishes the more general point relevant to this thesis's own basic distillation experiment: plain, spatially-unaware logit- or feature-matching, adequate for image classification, requires this kind of foreground/background- and attention-aware treatment once it is applied to a detection head, and a naive transplant of the classification recipe should not be expected to perform as well on a detector without it.

Beyond this foundational material, a substantial further literature addresses problems this thesis's basic comparison does not engage with: the \emph{capacity gap}, whereby distillation quality degrades -- and can invert, so that a smaller or weaker teacher produces a better student than a larger, more accurate one -- once the disparity between teacher and student capacity grows too large, together with the associated \emph{teacher-assistant} remedy of interposing one or more intermediate-capacity networks between teacher and student to bridge that gap \cite{mirzadehImprovedKnowledgeDistillation2020}; distillation schemes built specifically around relation-based knowledge; and cross-domain or cross-architecture distillation, relevant when teacher and student differ in label space or backbone family. None of these were investigated in this thesis; the single basic distillation comparison reported in \Cref{sec:results_kd_and_embedded} is interpreted against the foundational material above alone.

%%%%%%%%%%%%%%%%%%%%%%%%%%%%%%%%%%%%%%%%%%%%%%%%%
\subsection{Embedded Inference Feasibility}\label{sec:lit_embedded_inference}

Three design principles recur across the lightweight architectures considered as distillation-student candidates and surveyed among the YOLO family's nano-scale variants (\Cref{sec:lit_yolo_family}). \emph{Depthwise separable convolutions}, introduced for mobile vision by \textit{MobileNet} \cite{howardMobileNetsEfficientConvolutional2017}, factorize a standard convolution into a per-channel spatial filtering step and a $1\times1$ pointwise combination step, decoupling the interaction between kernel size and channel count that otherwise dominates a standard convolution's computational cost. \emph{Anchor-free detection heads} remove the need to hand-design and re-tune anchor boxes per dataset; \textit{PP-PicoDet} \cite{yuPPPicoDetBetterRealTime2021} pairs an anchor-free head with a lightweight neck and an enhanced ShuffleNetV2-derived backbone to reach $30.6\%$ mAP at $0.99$M parameters, outperforming the similarly-sized anchor-based \textit{YOLOX-Nano} at lower measured mobile-CPU latency. \emph{Channel/filter pruning} \cite{liPruningFiltersEfficient2017} instead starts from a trained, larger network and removes whole filters identified as contributing little to output accuracy, which -- unlike unstructured weight-level pruning -- yields a smaller but still dense network that runs on ordinary convolution kernels without requiring sparse-matrix hardware support. These three principles are the conceptual background for why nano-scale detectors were the candidates of interest as a distillation student (\Cref{sec:lit_yolo_family}, \Cref{sec:model_selection}); none of them were themselves implemented or modified in this thesis.

Standardizing how on-device inference is measured is itself an active concern in the literature and in industry practice. \textit{MLPerf Inference} \cite{reddiMLPerfInferenceBenchmark2020} established an architecture-neutral methodology for comparing inference latency and throughput across heterogeneous ML hardware and software stacks; \textit{MLPerf Tiny} \cite{banburyMLPerfTinyBenchmark2021} adapts this methodology to microcontroller- and DSP-class devices, reporting accuracy, latency, and energy together and standardizing on a common reference inference stack so that cross-device comparisons are not confounded by inconsistent runtime choices. Consistent with this, \textit{MLPerf Tiny}'s own measurement protocol deliberately scopes the timed window to model inference alone, explicitly excluding any pre- or post-processing performed around it -- a methodological control this thesis's own runtime check on the \textit{AX Visio} hardware (\Cref{sec:results_ax_visio_benchmark}) also needs to observe, alongside the related, standard practical concern of excluding one-time effects (initial model loading, cache and frequency-scaling warm-up) from a measurement taken as a single or small number of runs. \textit{TensorFlow Lite Micro} \cite{davidTensorFlowLiteMicro2021} illustrates why the choice of export format and runtime matters for such a benchmark's validity in the first place: it argues for a portable, interpreter-based deployment framework over a bespoke per-device compiled pipeline precisely because the embedded ecosystem is fragmented across instruction sets, absent OS/file-system/dynamic-memory support, and vendor-specific kernel optimizations, and it separately documents that its interpreter's behavior under multi-threaded or multi-core execution is a first-class design concern rather than an afterthought -- underscoring that thread and core-affinity choices on a multi-core ARM SoC (such as the \textit{AX Visio}'s or the \textit{Raspberry Pi 5}'s) are a further variable this thesis's benchmark protocol needs to hold fixed across any repeated measurement, alongside the export format (ONNX, TFLite, NCNN, or a vendor-specific format such as Qualcomm's SNPE \texttt{.dlc}) used to produce the deployed model.

On the numerical-precision side, quantizing a trained network's weights and activations from 32-bit floating point to 8-bit integers is well established as a first, comparatively cheap optimization step, and is directly relevant to the \textit{Hexagon 685} DSP because both foundational quantization papers in this literature explicitly target Qualcomm's fixed-point DSP hardware as their deployment case. \citeauthor{krishnamoorthiQuantizingDeepConvolutional2018}'s whitepaper on post-training quantization (PTQ) \cite{krishnamoorthiQuantizingDeepConvolutional2018} reports that per-channel weight / per-layer activation 8-bit quantization keeps classification accuracy within $2\%$ of the floating-point baseline across a range of CNN architectures, and benchmarks the resulting speedup on \enquote{Qualcomm QDSPs with HVX} -- the same class of fixed-point SIMD DSP as the Hexagon 685 -- at up to $10\times$ over floating-point CPU inference. \citeauthor{jacobQuantizationTrainingNeural2018} \cite{jacobQuantizationTrainingNeural2018} extend this to quantization-aware training (QAT), in which quantization is simulated during training via \enquote{fake quantization} nodes inserted into the computation graph so that weights and activation ranges adapt to quantization noise before it is applied for inference; this paper explicitly frames \enquote{integer-arithmetic-only hardware such as the Qualcomm Hexagon} as its target deployment setting, and reports that QAT narrows the residual accuracy gap to the floating-point baseline further than PTQ alone. On the \textit{QCS605} specifically, the vendor toolchain implementing this PTQ/QAT distinction is the Snapdragon Neural Processing Engine (SNPE) together with the AI Model Efficiency Toolkit (AIMET) \cite{QualcommAimet2026} -- vendor documentation rather than peer-reviewed literature, cited here only to identify the toolchain rather than as an independent technical source. This thesis surveys quantization for this context but does not implement PTQ, QAT, or any other quantization or compression technique: the result reported in \Cref{sec:results_ax_visio_benchmark} is a single runtime measurement of the unoptimized, full-precision model, and the quantization pass sketched in the original research plan was not carried out.

Development and the majority of benchmarking in this thesis targeted the \textit{Raspberry Pi 5} as a stand-in for the \textit{QCS605} rather than the \textit{QCS605} itself, primarily for software-maturity reasons: the \textit{QCS605}'s Hexagon 685 DSP and Adreno 615 GPU are tied to an aging, 2018-era vendor software stack whose limited and inconsistently documented operator support makes routine model iteration as much a driver-debugging exercise as a modeling one. The \textit{Raspberry Pi 5}'s ARM Cortex-A76 CPU cores, by contrast, are estimated to run on the order of $60\%$ faster than the \textit{QCS605}'s CPU cores under comparable workloads, which gives development a performance \emph{ceiling} rather than a performance floor: a model shown to be fast enough on the \textit{Raspberry Pi 5}, after discounting for this margin, is expected to be fast enough on the \textit{QCS605}, whereas the converse inference -- that a model too slow on the \textit{Raspberry Pi 5} might still be acceptable on the \textit{QCS605} -- does not hold with the same confidence. The single runtime check on the \textit{AX Visio} hardware reported in \Cref{sec:results_ax_visio_benchmark} accordingly functions as a validation of this proxy relationship on real target-adjacent hardware, rather than as a replacement for the \textit{Raspberry Pi 5} as the thesis's primary development and benchmarking platform.
